\section{GDT 描述符的 8 字节布局}

这是 x86 最"反直觉"的数据结构之一——由于历史兼容性（286→386），
字段被拆成不连续的碎片散布在 8 个字节中。

\subsection{内存中的字节布局}

\begin{figure}[H]
  \centering
  % figures/ch03/fig03-gdt-layout.tex — auto-extracted from chapter source
\begin{tikzpicture}[node distance=0cm]
  % 8 bytes, each a row
  \node[memcell, minimum width=6cm, fill=green!10] (b0) {limit\_low [7:0]};
  \node[memaddr, left=0.3cm of b0] {+0};

  \node[memcell, minimum width=6cm, fill=green!10, above=0cm of b0] (b1) {limit\_low [15:8]};
  \node[memaddr, left=0.3cm of b1] {+1};

  \node[memcell, minimum width=6cm, fill=blue!10, above=0cm of b1] (b2) {base\_low [7:0]};
  \node[memaddr, left=0.3cm of b2] {+2};

  \node[memcell, minimum width=6cm, fill=blue!10, above=0cm of b2] (b3) {base\_low [15:8]};
  \node[memaddr, left=0.3cm of b3] {+3};

  \node[memcell, minimum width=6cm, fill=blue!10, above=0cm of b3] (b4) {base\_mid [23:16]};
  \node[memaddr, left=0.3cm of b4] {+4};

  \node[memcell, minimum width=6cm, fill=red!10, above=0cm of b4] (b5) {access byte};
  \node[memaddr, left=0.3cm of b5] {+5};

  \node[memcell, minimum width=6cm, fill=yellow!15, above=0cm of b5] (b6) {flags[7:4] | limit\_high[19:16]};
  \node[memaddr, left=0.3cm of b6] {+6};

  \node[memcell, minimum width=6cm, fill=blue!10, above=0cm of b6] (b7) {base\_high [31:24]};
  \node[memaddr, left=0.3cm of b7] {+7};

  % Braces on the right
  \draw[decorate, decoration={brace, amplitude=8pt}]
    ([xshift=0.3cm]b1.north east) -- ([xshift=0.3cm]b0.south east)
    node[midway, right=10pt, font=\small, text=green!50!black] {limit[0:15]};

  \draw[decorate, decoration={brace, amplitude=8pt}]
    ([xshift=0.3cm]b4.north east) -- ([xshift=0.3cm]b2.south east)
    node[midway, right=10pt, font=\small, text=blue!50!black] {base[0:23]};

  \draw[decorate, decoration={brace, amplitude=8pt}]
    ([xshift=0.3cm]b7.north east) -- ([xshift=0.3cm]b7.south east)
    node[midway, right=10pt, font=\small, text=blue!50!black] {base[24:31]};
\end{tikzpicture}

  \caption{GDT 描述符的 8 字节内存布局（地址从下到上递增）}
  \label{fig:gdt-layout}
\end{figure}

\subsection{对应的 C 结构体}

\codefile{src/include/arch/gdt.h}
\begin{lstlisting}[style=cstyle]
typedef struct __attribute__((packed)) {
    uint16_t limit_low;   // 字节 0-1: limit 的低 16 位 [0:15]
    uint16_t base_low;    // 字节 2-3: base 的低 16 位 [0:15]
    uint8_t  base_mid;    // 字节 4:   base 的中间 8 位 [16:23]
    uint8_t  access;      // 字节 5:   访问字节 (P/DPL/S/Type)
    uint8_t  gran;        // 字节 6:   高4位=flags, 低4位=limit[16:19]
    uint8_t  base_high;   // 字节 7:   base 的最高 8 位 [24:31]
} gdt_entry_t;
\end{lstlisting}

\begin{warnbox}
\texttt{\_\_attribute\_\_((packed))} 是\textbf{必须的}。没有它，编译器可能在
\texttt{base\_mid} 后面插入 padding 字节来对齐 \texttt{access}，导致整个结构体
从 8 字节膨胀到 12 字节——CPU 会把 padding 当作描述符的一部分，结果完全错乱。
\end{warnbox}

\begin{thinkbox}
为什么 base 被拆成三段（low/mid/high）而非一个连续的 32 位字段？

因为 286 时代 base 只有 24 位（3 字节），386 扩展到 32 位时在末尾加了 1 字节。
Intel 必须向后兼容 286 的描述符格式，所以新增的 8 位只能塞到字节 7。
这就是为什么 base 是"碎片化"的——历史包袱。
\end{thinkbox}

\subsection{重组逻辑：\texttt{gdt\_set\_entry}}

因为字段是碎片化的，填充描述符需要位操作来拆分和重组：

\codefile{src/kernel/arch/gdt.c}
\begin{lstlisting}[style=cstyle]
static void gdt_set_entry(int idx, uint32_t base, uint32_t limit,
                          uint8_t access, uint8_t flags) {
    // --- limit 拆分：低 16 位放 limit_low，高 4 位放 gran 的低半 ---
    gdt[idx].limit_low = (uint16_t)(limit & 0xFFFF);
    //   limit = 0x000FFFFF
    //   limit & 0xFFFF = 0xFFFF → limit_low

    // --- base 拆分：3 段 ---
    gdt[idx].base_low  = (uint16_t)(base & 0xFFFF);
    //   base = 0x00000000
    //   base & 0xFFFF = 0x0000 → base_low

    gdt[idx].base_mid  = (uint8_t)((base >> 16) & 0xFF);
    //   (base >> 16) & 0xFF = 0x00 → base_mid

    gdt[idx].base_high = (uint8_t)((base >> 24) & 0xFF);
    //   (base >> 24) & 0xFF = 0x00 → base_high

    // --- access 直接赋值（完整的 8 位） ---
    gdt[idx].access = access;

    // --- gran 字节：高 4 位是 flags，低 4 位是 limit 的 [16:19] ---
    gdt[idx].gran = (uint8_t)((flags & 0xF0) | ((limit >> 16) & 0x0F));
    //   flags = 0xC0 → 高 4 位 = 0xC (= 1100b → G=1, D/B=1)
    //   (limit >> 16) & 0x0F = 0x0F → 低 4 位
    //   结果: 0xCF
}
\end{lstlisting}

对于我们的内核代码段（base=0, limit=\hex{FFFFF}, access=\hex{9A}, flags=\hex{C0}），
实际写入内存的 8 字节是：

\begin{figure}[H]
  \centering
  % figures/ch03/fig04-8.tex — auto-extracted from chapter source
\begin{tikzpicture}[node distance=0cm]
  \foreach \i/\val/\desc in {
    0/\texttt{0xFF}/limit\_low[7:0],
    1/\texttt{0xFF}/limit\_low[15:8],
    2/\texttt{0x00}/base\_low[7:0],
    3/\texttt{0x00}/base\_low[15:8],
    4/\texttt{0x00}/base\_mid,
    5/\texttt{0x9A}/access,
    6/\texttt{0xCF}/flags|limit\_hi,
    7/\texttt{0x00}/base\_high} {
    \node[memcell, minimum width=1.5cm, minimum height=0.7cm, fill=white]
      (byte\i) at (\i*1.6, 0) {\val};
    \node[font=\tiny\ttfamily, above=0.1cm of byte\i] {+\i};
    \node[font=\tiny, below=0.05cm of byte\i, text=gray] {\desc};
  }
\end{tikzpicture}

  \caption{内核代码段描述符在内存中的实际 8 字节}
\end{figure}

% ============================================================================

